\newpage{\ }
\thispagestyle{empty}

\begin{spacing}{1.0} % Espaciado 1.0 solo aquí

\chapter*{Summary}
\addcontentsline{toc}{chapter}{Summary}

\vspace{.5cm}

Northwestern Patagonian is a fire-prone landscape. With climate change, this region is expected to experience increased temperature and decreased precipitation, which would significantly enhance fire activity. In this context, advancing our understanding of the various factors that control fire is crucial to anticipate changes in the fire regime. Additionally, it is necessary to develop models that represent and simulate these processes to project future fire regimes under changing environmental conditions. In this thesis, we study how physical factors (climate and topography), vegetation, and human presence influence fire in this region. We developed models to simulate fire regimes and project fire activity under future climate scenarios. 

Before developing simulation models, we examined the physical and biotic controls of fire at two scales. At the site scale, we evaluated how the contrasting microclimate of adjacent plant communities regulates fuel moisture, contributing to differences in flammability. At the landscape scale, we analyze how climate and topography regulate fire activity by interacting with vegetation type.

In northwestern Patagonia, two vegetation types dominate the landscape: shrublands, which are relatively flammable and dominated by resprouting species, and forests, which are less flammable and less resilient to fire. Their differences in flammability are largely explained by their contrasting structure: shrublands have a higher amount and continuity of fine fuels. By comparing adjacent forests and shrublands, we found that forests also have a cooler and more humid microclimate than shrublands, where dead fuel moisture is higher. Thus, their lower flammability can be explained not only by their lower fine fuel amount and continuity but also by higher dead fuel moisture.

The lower flammability of forests is also reflected in a lower fire activity when analyzed at the landscape scale. However, at this scale, other factors can confound the effect of vegetation. For example, forests tend to occupy areas with higher precipitation and elevation (lower temperature), while shrublands are found closer to human settlements, where ignition frequency is higher. To properly model the fire regime while controlling for relevant variables, we analyzed how topography and mean annual precipitation regulate fire activity while interacting with vegetation type. We mapped all fires that occurred in a 24-year period (1999-2022) in northwestern Patagonia (39° to 44.5° S) and assessed how the probability of grid cells burning at least once in the study period varied as a function of environmental variables. We found that fire activity decreases with elevation and precipitation, but the effect of precipitation is largely due to how vegetation types with contrasting structure are distributed along the gradient: vegetation types with higher fine fuel amount and continuity and lower intrinsic fuel moisture are found in low and intermediate precipitation areas. Additionally, we detected a clear unimodal effect of productivity, both within and among vegetation types.

Building on our understanding of how different factors control fire activity, we developed a fire regime simulator composed of three spatially explicit fire models: ignition, escape, and spread. The models were defined in discrete time and space, with a spatial resolution of 30 m and a temporal resolution of 14 days. Ignitions can be caused by lightning or human activity, and once they occur, they may completely burn the original cell and escape from it. Given the escape, the spread model regulates how fire propagates cell by cell. All models include the effects of atmospheric conditions through the Fire Weather Index, vegetation (vegetation type and productivity), and topography. Additionally, the human ignition and escape models consider the distance to the nearest settlement and road, while the spread model incorporates wind effects. We estimated the ignition and escape models using fire records from Nahuel Huapi National Park, while the spread model was estimated using mapped fire polygons from the previous analysis. In all three models, we found a clear effect of fire weather, with the Fire Weather Index increasing the probability of ignition, escape, and spread. We found that, although human ignitions are more frequent than lightning ignitions, the latter have a higher probability of escape because they occur far from roads, making early suppression more difficult. Using these models, we projected fire activity in Nahuel Huapi National Park for the periods 2040-2049 and 2090-2099, considering four contrasting climate scenarios. The simulations indicate that for 2040-2049, the annual burn probability will increase by a factor of 2.07 to 2.61, while for 2090-2099, increases are expected by factors ranging from 1.36 to 31.17, depending on the climatic scenario.

Our study highlights the importance of implementing effective fire management policies to mitigate its effects on societies and ecosystems. In particular, with increasing fire activity, the Andean-Patagonian forests are expected to decline, at the expense of the expansion of shrublands and grasslands. These transitions, in turn, would further promote fire occurrence.

\noindent \textbf{Keywords:}  \textit{Burn probability, models, fire spread, vegetation, topography.}

\end{spacing}